\chapter{Growth Empirics and Convergence}\label{ch:b3:growth-empirics}

The Year 1 volume built the Solow model and the Year 2 volume the models of growth
through ideas. This chapter takes them to the data of some 180 countries over sixty
years. Do poor countries catch up? Which national characteristics go with fast growth,
and how much can a regression across countries say? Why do some economies take off
and others stall? And how much of growth comes from workers leaving farms? The data
answer some of these questions well and others hardly at all; the chapter separates
the two.

\begin{openingchart}{Annual growth of real GDP per person by country, decade by
decade (the 2010s are 2010--2019). Each dot is a country; the diamonds are the
medians. Values beyond $-8$ and $+12$~\% are drawn at those limits. Growth is measured
with the national-accounts series, which the source recommends for changes over
time.}{Penn World Table 11.0, rgdpna and pop}
\begin{tikzpicture}
\begin{axis}[omscatter, width=0.72\linewidth, height=0.44\linewidth, xmin=-0.6,
  xmax=5.6, ymin=-8.5, ymax=12.5, xtick={0,...,5},
  xticklabels={1960s, 1970s, 1980s, 1990s, 2000s, 2010s},
  ylabel={Percent a year}, mark size=0.8pt,
  extra y ticks={0}, extra y tick labels={}, extra y tick style={grid=major}]
  \addplot table[x=pos, y=growth] {figdata/bachelor-3/31-growth-empirics/out/decades.csv};
  \addplot[only marks, mark=diamond*, mark size=3pt, omThm] table[x=pos, y=median] {figdata/bachelor-3/31-growth-empirics/out/decades_median.csv};
\end{axis}
\end{tikzpicture}
\end{openingchart}

% figdata: decades_stats
The median country grew 2.9~\% a year in the 1960s, 1.2~\% in the 1980s and 2.4~\% in the 2000s. In
every decade the spread is wide: the middle half of countries spans 2.6 to 3.8 points,
and in the 1980s 37~\% of countries saw income per person fall. A country's place in one
decade is a weak guide to its place in the next: the correlation of growth between
successive decades is 0.45 from the 1960s to the 1970s, 0.33 from the 1980s to the 1990s,
$-0.08$ from the 1990s to the 2000s and 0.32 from the 2000s to the 2010s.

\begin{recall}
The Solow model, its steady state and the speed of convergence, and growth
accounting are the Year 1 volume's, as is regression to the mean; the Ramsey model and
growth through ideas are the Year 2 volume's. Panel data and fixed effects are
\cref{ch:b3:panel-data}'s, specification search \cref{ch:b3:credible-empirical-work}'s.
\end{recall}

\begin{notation}[Symbols of this chapter]\label{not:b3:growth-empirics:symbols}
$y_{it}$ is real GDP per person; $g_i$ the average annual log growth over $T$ years; $\beta$ the slope of
growth on $\ln y_{i0}$; $\lambda$ the speed of convergence; $\sigma_t$ the cross-country standard deviation of $\ln y_{it}$; $\alpha$
the capital share.
\end{notation}

\section{Convergence}\label{sec:b3:growth-empirics:convergence}

\begin{definition}[Beta convergence; sigma convergence]\label{def:b3:growth-empirics:convergence}
\index{beta convergence}\index{sigma convergence}
Countries show \emph{beta convergence} if, in a regression of growth on initial log
income, the slope $\beta$ is negative: the initially poor grow faster. They show \emph{sigma
convergence} if the dispersion of log income $\sigma_t$ falls over time.
\end{definition}

\begin{proposition}[Reading the slope as a speed]\label{prop:b3:growth-empirics:speed}
Near its steady state $y^*$, the Solow model gives $\ln y_T - \ln y_0 = -(1 - e^{-\lambda T})(\ln y_0 - \ln y^*)$, with $\lambda = (1 - \alpha)(n + g + \delta)$. A
regression of average growth on $\ln y_0$ among countries with the same $y^*$ therefore has slope
$\beta = -(1 - e^{-\lambda T})/T$, and $\lambda = -\ln(1 + \beta T)/T$.
\end{proposition}

\begin{proof}
The Year 1 volume linearised the capital accumulation equation around the steady
state: $d\ln y/dt = -\lambda(\ln y - \ln y^*)$. Integrating from 0 to $T$ gives $\ln y_T - \ln y^* = e^{-\lambda T}(\ln y_0 - \ln y^*)$; subtracting $\ln y_0 - \ln y^*$ from both
sides gives the first equation. Dividing by $T$, growth is linear in $\ln y_0$ with slope $-(1 - e^{-\lambda T})/T$;
solving for $\lambda$ gives the last expression.
\end{proof}

With $\alpha = 1/3$ and $n + g + \delta$ near 6~\%, the model predicts $\lambda$ near 4~\% a year.

\begin{proposition}[Beta does not imply sigma]\label{prop:b3:growth-empirics:galton}
\omterm{def:b3:growth-empirics:convergence}{Beta convergence} is necessary for \omterm{def:b3:growth-empirics:convergence}{sigma convergence} but not sufficient: with shocks,
the poor can grow faster on average while the dispersion of incomes stays constant.
\end{proposition}

\begin{proof}
Let $\ln y_{i1} = a + (1 + b)\ln y_{i0} + e_i$, with shocks $e_i$ independent of $\ln y_{i0}$ and variance $s^2$, so that the slope of growth on
$\ln y_{i0}$ is $b$. Then $\sigma_1^2 = (1 + b)^2\sigma_0^2 + s^2$. If $b \ge 0$, $\sigma_1 \ge \sigma_0$: no \omterm{def:b3:growth-empirics:convergence}{sigma convergence} without \omterm{def:b3:growth-empirics:convergence}{beta convergence}.
If $b < 0$ but $s^2 = (1 - (1 + b)^2)\sigma_0^2$, then $\sigma_1 = \sigma_0$: \omterm{def:b3:growth-empirics:convergence}{beta convergence} without \omterm{def:b3:growth-empirics:convergence}{sigma convergence}. This is
regression to the mean.
\end{proof}

\begin{omchart}{Annual growth of real GDP per person, 1960--2019, against GDP per person
in 1960 (logarithmic scale). Filled dots: Western Europe and its offshoots (Maddison's
regions), a group fixed by geography; with least-squares lines for all countries and
for the group. Growth from the national-accounts series, the 1960 level from the
expenditure-side series, which compares countries in one year.}{Penn World Table 11.0,
rgdpna, rgdpe and pop; Maddison Project Database 2023, regions}
\begin{tikzpicture}
\begin{axis}[omscatter, width=0.62\linewidth, height=0.42\linewidth,
  xlabel={GDP per person in 1960, dollars}, ylabel={Growth 1960--2019, \% a year},
  xtick={2.5,3,3.5,4,4.5}, xticklabels={300, 1{,}000, 3{,}000, 10{,}000, 30{,}000},
  scatter/classes={False={mark=o, draw=black!50}, True={mark=*, fill=omDef, draw=omDef}},
  scatter, point meta=explicit symbolic]
  \addplot table[x=lx, y=growth, meta=rich] {figdata/bachelor-3/31-growth-empirics/out/convergence.csv};
  \addplot[black!60, thick, no marks, sharp plot] table[x=lx, y=growth] {figdata/bachelor-3/31-growth-empirics/out/convergence_fit_all.csv};
  \addplot[omThm, thick, no marks, sharp plot] table[x=lx, y=growth] {figdata/bachelor-3/31-growth-empirics/out/convergence_fit_rich.csv};
\end{axis}
\end{tikzpicture}
\end{omchart}
% figdata: convergence_stats

Across all 111 countries with data for 1960 there is no absolute convergence: the slope is
$-0.0011$, with a standard error of 0.0014. Among the 23 countries of Western Europe and its offshoots
the slope is $-0.0108$ (0.0018), a speed of convergence of 1.7~\% a year --- under half what the
Solow model predicts with standard parameters.

\begin{model}{The augmented Solow model in cross-country data}
\assumes Countries share one technology; saving, schooling and population growth
are given; output depends on physical and \omterm{def:b3:human-capital:human-capital}{human capital}.
\predicts Conditional convergence: holding the determinants of the steady state
fixed, poor countries grow faster, at a speed the model's parameters fix.
\tested The augmented model explains about 80~\% of the international variation in
income per person, and, holding population growth and capital accumulation
constant, countries converge at about the predicted rate. % ledger: b3ge:mrw
\breaks Physical and \omterm{def:b3:human-capital:human-capital}{human capital} explain only part of output per worker; the
Solow residual varies greatly across countries, with institutions and policies
behind it. % ledger: b3ge:hj
With country effects in a panel, which let production functions differ, the
estimates differ significantly from those of a single cross-section. % ledger: b3ge:islam
\end{model}

The panel estimates are not simply better. As \cref{ch:b3:panel-data} showed, a
lagged dependent variable with a fixed effect is biased in short panels, and short
periods mix convergence with the business cycle.

\section{Growth regressions and their fragility}\label{sec:b3:growth-empirics:regressions}

\begin{definition}[Growth regression; extreme-bounds analysis]\label{def:b3:growth-empirics:growth-regression}
\index{growth regression}\index{extreme-bounds analysis}
A \emph{growth regression} relates countries' average growth to their initial income
and other characteristics. \emph{Extreme-bounds analysis} re-estimates the coefficient
of a variable with every combination of other candidate controls and calls it robust
if the lowest and highest estimates (less and plus two standard errors) have the same
sign.
\end{definition}

\begin{record}[Fragile determinants]
With a hundred or so countries and dozens of candidate variables, results depend on
which controls enter. Extreme bounds found almost all results fragile; only the
share of investment in GDP was robustly correlated with growth, and investment with
trade. % ledger: b3ge:lr
Looking instead at the whole distribution of a variable's estimates across millions
of regressions, a substantial number of variables are strongly related to growth. % ledger: b3ge:sim
The two conclusions use the same data and differ in what they call robust.
Exercise 11 shows why extreme bounds is a harsh test: a variable with a true effect can
fail it when its candidate controls are correlated with it.
\end{record}

\begin{measure}[Which vintage of the data?]
Penn World Table estimates of growth change from one version to the next, though the
underlying data and methods are almost the same; the variation is systematic and comes
from the method. Studies that use growth over long periods are robust to the revisions;
studies of annual growth are much less so. % ledger: b3ge:jlps
This book uses version 11.0; an exercise run on another version would give somewhat
different numbers, and long averages should be preferred. The choice of series within
a version matters too: the expenditure-side series compares countries' living
standards in one year, but it also moves with relative prices and the \omterm{def:b3:ricardian-trade:terms-of-trade}{terms of trade},
so this chapter measures growth with the national-accounts series, as the table's
documentation advises; with the expenditure-side series the correlation of growth
between the 2000s and the 2010s is $-0.09$ instead of 0.32.
\end{measure}

Growth is also hard to explain because it does not persist. Country characteristics
such as schooling are stable from decade to decade, with correlations of 0.6 to 0.9, while
growth rates correlate at only 0.1 to 0.3; shocks, especially to the \omterm{def:b3:ricardian-trade:terms-of-trade}{terms of trade}, account for
much of the variation. % ledger: b3ge:ekps

\section{Accelerations, miracles and collapses}\label{sec:b3:growth-empirics:accelerations}

\begin{definition}[Growth acceleration; growth collapse]\label{def:b3:growth-empirics:acceleration}
\index{growth acceleration}\index{growth collapse}
A \emph{growth acceleration} is a rise in a country's growth rate, to a rapid pace,
sustained for at least eight years. A \emph{growth collapse} is a sustained fall of
income per person.
\end{definition}

% ledger: b3ge:hpr
Since the 1950s there have been more than 80 accelerations. They go with rising investment
and trade and with real depreciation; changes of political regime predict them;
those that follow external shocks fizzle out, while those that follow economic
reform tend to last. But most accelerations are unrelated to the usual determinants,
and most reforms produce none: turning points are hard to predict.

The filter used here to date one is the same in spirit as the study's: year $t$
qualifies when mean growth over the next eight years is at least 3.5~\%, at least 2
points above the mean of the past five, and income at the end is a new peak. Several
consecutive years usually qualify for one episode; of those, the filter keeps the year
with the largest rise in growth, which the study's own choice of the best-fitting
break resembles.

\begin{omchart}{GDP per person, 1900--2022, five economies (logarithmic scale, 2011
dollars).}{Maddison Project Database 2023}
\begin{tikzpicture}
\begin{semilogyaxis}[omtime, width=0.7\linewidth, height=0.44\linewidth, xmin=1900,
  xmax=2022, ymin=700, ymax=80000, ylabel={Dollars},
  yticklabels={1{,}000, 10{,}000}, ytick={1000,10000},
  legend pos=north west, legend style={font=\footnotesize}, legend columns=2]
  \addplot[black, thick] table[x=year, y=y] {figdata/bachelor-3/31-growth-empirics/out/paths_USA.csv};
  \addplot[omDef, thick] table[x=year, y=y] {figdata/bachelor-3/31-growth-empirics/out/paths_ARG.csv};
  \addplot[omThm, thick] table[x=year, y=y] {figdata/bachelor-3/31-growth-empirics/out/paths_JPN.csv};
  \addplot[omThm, dashed, thick] table[x=year, y=y] {figdata/bachelor-3/31-growth-empirics/out/paths_KOR.csv};
  \addplot[omDef, dashed, thick] table[x=year, y=y] {figdata/bachelor-3/31-growth-empirics/out/paths_VEN.csv};
  \legend{United States, Argentina, Japan, South Korea, Venezuela}
\end{semilogyaxis}
\end{tikzpicture}
\end{omchart}
% figdata: paths_points

\begin{record}[Growth that stops: Argentina and Venezuela]
In 1913 Argentina's income per person was 6{,}052 dollars of 2011, 60~\% of the U.S. level; in 2022 it was 18{,}292,
31~\% of it. Venezuela was richer than Argentina in 1950, at 8{,}462 dollars, and reached 16{,}270 in 1980; in
2022 it was back to 5{,}267, below its 1950 level. Japan went from 3{,}062 in 1950 to 38{,}269 in 2022, and South
Korea from 998 to 41{,}321. % figdata: paths_points
\end{record}

\begin{record}[The East Asian growth miracle]
Over the two and a half decades to the early 1990s, output per person grew 6 to 7~\% a
year in Hong Kong,
Singapore, South Korea and Taiwan. Participation, schooling and investment rose
quickly; outside Hong Kong, non-farm employment grew one and a half to two times as
fast as the working population. Output per effective worker outside farming grew
only 3 to 4~\% a year, and once the doubling to quadrupling of investment rates is
allowed for, total factor productivity grew at rates within those of other economies
over equally long periods. The growth was extraordinary; its productivity growth was
not. % ledger: b3ge:young
\end{record}

\section{The move out of agriculture}\label{sec:b3:growth-empirics:structural}

\begin{definition}[Structural transformation; agricultural productivity gap]\label{def:b3:growth-empirics:structural}
\index{structural transformation}\index{agricultural productivity gap}
\emph{Structural transformation} is the shift of workers and output from agriculture
to industry and services as income grows. The \emph{agricultural productivity gap} is
the ratio of value added per worker outside agriculture to value added per worker in
it.
\end{definition}

\begin{omchart}{Employment in agriculture against GDP per person (PPP, constant 2021
dollars, logarithmic scale), 2019, 177 countries.}{World Bank, World Development
Indicators, SL.AGR.EMPL.ZS (modelled ILO estimates) and NY.GDP.PCAP.PP.KD}
\begin{tikzpicture}
\begin{axis}[omscatter, width=0.6\linewidth, height=0.4\linewidth,
  xlabel={GDP per person, dollars}, ylabel={Agriculture, \% of employment},
  xtick={3,4,5}, xticklabels={1{,}000, 10{,}000, 100{,}000}, ymin=0, ymax=90]
  \addplot table[x=lx, y=agri] {figdata/bachelor-3/31-growth-empirics/out/agri.csv};
\end{axis}
\end{tikzpicture}
\end{omchart}
% figdata: agri_stats

In the poorest fifth of countries the median share of workers in agriculture is 45~\%; in
the richest fifth it is 2~\%. Value added per worker is much higher outside agriculture than
in it, especially in poor countries. Taken at face value the gap means that moving
workers out of farming raises output; after correcting for hours and schooling, and
with output measured from household surveys, a large gap remains. % ledger: b3ge:glw
How much of growth the move explains is a decomposition of the kind built in
\cref{ch:b3:platforms-and-innovation}, across sectors instead of firms; the weekend
problem performs one.

\section{Exercises}\label{sec:b3:growth-empirics:exercises}

\begin{exercise}[$\star$]\label{exo:b3:growth-empirics:1}
In the opening chart, in which decade was the median lowest, and what share of
countries saw income per person fall in it?
\end{exercise}

\begin{exercise}[$\star$]\label{exo:b3:growth-empirics:2}
A regression over 40 years gives $\beta = -0.015$. Find the speed of convergence.
\end{exercise}

\begin{exercise}[$\star$]\label{exo:b3:growth-empirics:3}
How long does income take to double at 2~\% and at 7~\% a year?
\end{exercise}

\begin{exercise}[$\star$]\label{exo:b3:growth-empirics:4}
What data would contradict conditional convergence?
\end{exercise}

\begin{exercise}[$\star\star$]\label{exo:b3:growth-empirics:5}
Estimate absolute convergence from 1960 to 2019 for all countries, for the 25 richest in
1960 and for the 25 richest in 2019. Why does the last sample show convergence?
\end{exercise}

\begin{exercise}[$\star\star$]\label{exo:b3:growth-empirics:6}
Compute the correlation of growth between successive decades, and of log income at the
start of successive decades.
\end{exercise}

\begin{exercise}[$\star\star$]\label{exo:b3:growth-empirics:7}
In the augmented Solow model with capital shares $\alpha = \phi = 1/3$ for physical and \omterm{def:b3:human-capital:human-capital}{human
capital}, steady-state income is proportional to
$(s_k^{\alpha}s_h^{\phi}/(n + g + \delta)^{\alpha + \phi})^{1/(1 - \alpha - \phi)}$. Country A saves 24~\% and invests 12~\% in schooling; country B saves 12~\% and
invests 6~\%, with the same $n + g + \delta$. Find the ratio of their steady-state incomes.
\end{exercise}

\begin{exercise}[$\star\star$]\label{exo:b3:growth-empirics:8}
Growth is 1~\% a year for five years, then 5~\% for eight years. Does the chapter's
acceleration filter (mean growth over the next eight years at least 3.5~\%, at least 2 points above
the past five, and a new peak of income) date an acceleration?
\end{exercise}

\begin{exercise}[$\star\star$]\label{exo:b3:growth-empirics:9}
The dispersion of log incomes across countries rose from 1960 to 2000, while a regression
shows the poor growing faster in a sample. How is that possible?
\end{exercise}

\begin{exercise}[$\star\star\star$]\label{exo:b3:growth-empirics:10}
Derive the speed of convergence $\lambda = (1 - \alpha)(n + g + \delta)$ from the Solow model with $k$ per effective worker.
\end{exercise}

\begin{exercise}[$\star\star\star$]\label{exo:b3:growth-empirics:11}
On an invented table of 80 countries, run \omterm{def:b3:growth-empirics:growth-regression}{extreme-bounds analysis} for seven candidate
variables. Which are robust, and which variable with a true effect fails?
\end{exercise}

\begin{exercise}[$\star\star\star$]\label{exo:b3:growth-empirics:12}
Farming employs 40~\% of workers at output per worker 1, the rest at 3. If 10~\% of all workers move
out of farming at unchanged sector productivities, by how much does output per worker
rise? What if the true output per worker outside farming is 1.5, the gap halved?
\end{exercise}

\section{Problem: Brenland's Fifty Years}\label{sec:b3:growth-empirics:problem}

\begin{problem}[{Weekend problem --- Brenland's fifty years}]\label{pb:b3:growth-empirics:fifty}
Brenland's real GDP per person for 2000--2025 comes from its accounts
(\code{load("accounts")}, real GDP over population). Before 2000 the series is extended
with annual growth of 0.5~\% in 1976--1981, 4.0~\% in 1982--1995 and 2.0~\% in 1996--2000. Twenty invented
peers have incomes in 1975 spread around 9{,}000 dalens and grow until 2025 at
$2.5 - 1.2\ln(y_{75}/9{,}000)$~\% a year, plus noise. % figdata: brenland_stats

\medskip\noindent\textbf{Part I --- fifty years of growth.}
\begin{enumerate}
\item Find income per person in 1975, 2000 and 2025.
\item Find average growth in 1975--2000 and in 2000--2025.
\item At the 2000--2025 rate, how long does income take to double?
\item Over 1975--2025, by what factor did income grow?
\item Why do two decades of 4~\% growth matter more than they look?
\end{enumerate}

\medskip\noindent\textbf{Part II --- Brenland among its peers.}
\begin{enumerate}[resume]
\item Regress the peers' and Brenland's growth on the log of their 1975 incomes. Find $\beta$.
\item Find the implied speed of convergence.
\item What growth does the regression predict for Brenland, and how did it do?
\item Is this absolute or conditional convergence? What do the peers share?
\item Find the standard deviation of the peers' log incomes in 1975 and 2025. Explain the
change with \cref{prop:b3:growth-empirics:galton}.
\item Name one reason the estimated speed could be biased.
\end{enumerate}

\medskip\noindent\textbf{Part III --- an acceleration.}
\begin{enumerate}[resume]
\item Apply the chapter's acceleration filter to Brenland, keeping, of consecutive
qualifying years, the one with the largest rise. Which years qualify, and which does it date?
\item Find mean growth in the five years before and the eight years after.
\item Is the episode sustained?
\item Suppose Brenland reformed its trade rules in 1981. What would have to be checked
before calling the reform the cause?
\item Why would a regression of growth on policies over 1975--2025 miss the timing?
\end{enumerate}

\medskip\noindent\textbf{Part IV --- leaving the farms.} Farming had 20~\% of Brenland's workers in 1975
and 5~\% in 2025, with output per worker of 6 and 12 thousand dalens; the rest had 16 and 36 thousand.
\begin{enumerate}[resume]
\item Find output per worker in both years.
\item Find the \omterm{def:b3:growth-empirics:structural}{agricultural productivity gap} in both years.
\item Decompose the change into within, between and cross terms.
\item What share of the growth came from the move out of farming?
\item Why might this overstate the gain from reallocation?
\item How would hours worked and schooling change the measured gap?
\item What would Brenland's growth have been with the farm share frozen at 20~\%?
\item State the result: Brenland's convergence speed and the share of its growth from
the move out of farming.
\end{enumerate}
\end{problem}
